\begin{abstract}
QKF-Certifier checks source-bound certificates for a fragment of KnownBits abstract transformers. Its public release parses a restricted straight-line transfer language, replays local normalization steps, and checks the resulting coordinatewise mask function on nine one-bit abstract inputs. We prove that this finite check is necessary and sufficient for soundness at every positive width within the accepted fragment; equality of the nine output rows likewise characterizes optimality. The normalizer is separated from the replay checker, and a decreasing syntactic measure establishes termination of the release's rewrite system.

The construction is placed within the finite quotient framework of Papers I and II. For finite ground specifications of a separately compatible binary operation, we give a sound, complete relation-matrix closure with explicit quotient feedback, and a coverage theorem for sparse semilattice representations. These algebraic kernels have their own certificate contract and are distinguished from the public word-program checker. For the latter, the accepted coordinatewise language has exactly $2^{18}$ raw-mask semantic classes, $3^9$ everywhere-valid classes, and respectively $64$, $64$, and $16$ sound classes for AND, OR, and XOR. The sound classes form Boolean lattices under accumulation of correct precision facts; additional monotonicity reduces the first two counts to $26$.

The paper specifies the source semantics, all 29 rewrite-rule identifiers, trusted components, and the correspondence between claims, code, and tests. It is tied to QKF-Certifier v0.1.0a1 at an immutable commit. The supplement supplies reproducible finite checks and a version manifest so that later software and manuscript revisions can be compared without silently changing the scope of an earlier certificate.
\end{abstract}
\noindent\textbf{Keywords.} KnownBits; abstract transformers; finite quotients; congruence closure; proof certificates; term rewriting; width-independent verification.
